% Lista de pendientes de la web del lanzamiento de la cartilla.
% Compilar:  latexmk -pdf pendientes.tex
\documentclass[11pt,a4paper]{article}
\usepackage[spanish,es-noshorthands]{babel}
\usepackage[utf8]{inputenc}
\usepackage[T1]{fontenc}
\usepackage{lmodern}
\usepackage[margin=2.2cm]{geometry}
\usepackage[table]{xcolor}
\usepackage{tabularx,array,booktabs}
\usepackage{enumitem}
\usepackage{fontawesome5}
\usepackage[most]{tcolorbox}
\usepackage[hidelinks]{hyperref}

\definecolor{morado}{HTML}{2A1B7E}
\definecolor{amarillo}{HTML}{FFCF2E}
\definecolor{rojo}{HTML}{C4142C}
\definecolor{verde}{HTML}{1F9D55}
\definecolor{gris}{HTML}{5D5A73}

\newcommand{\alta}{\textcolor{rojo}{\faCircle}~Alta}
\newcommand{\media}{\textcolor{amarillo!80!black}{\faCircle}~Media}
\newcommand{\baja}{\textcolor{verde}{\faCircle}~Baja}
\newcommand{\casilla}{\faSquare[regular]}

\newcolumntype{Y}{>{\raggedright\arraybackslash}X}
\renewcommand{\arraystretch}{1.35}

% tabla de una sección: casilla | tarea | quién | prioridad
\newenvironment{pendientes}
  {\par\noindent\tabularx{\linewidth}{@{}c Y >{\raggedright\arraybackslash}p{2.6cm} p{1.7cm}@{}}
   \toprule \textbf{} & \textbf{Tarea} & \textbf{Quién} & \textbf{Prioridad}\\ \midrule}
  {\bottomrule\endtabularx\par\medskip}

\newcommand{\seccion}[2]{\section*{\textcolor{morado}{#1~#2}}}

\begin{document}

\begin{tcolorbox}[colback=morado,colframe=morado,coltext=white,arc=4mm,boxsep=6pt]
  {\Huge\bfseries\textcolor{amarillo}{Pendientes}}\\[4pt]
  {\Large Web del lanzamiento · \emph{Vivencia tecnológica: Sistemas mecánicos}}\\[2pt]
  GIVIE · ATA EPE Robótica Escolar · Corporación Escuela Pedagógica Experimental
\end{tcolorbox}

\noindent\begin{tabularx}{\linewidth}{@{}l Y@{}}
\faGlobe~\textbf{Web} & \url{https://sebrodriguezg.github.io/GIVIE-/}\\
\faGithub~\textbf{Repositorio} & \url{https://github.com/Sebrodriguezg/GIVIE-}\\
\faCalendar*~\textbf{Lanzamiento} & Sábado 7 de noviembre de 2026, 2:00--5:00 p.\,m., Corporación Escuela Pedagógica Experimental\\
\faClock~\textbf{Estado al} & 26 de octubre de 2026
\end{tabularx}

\begin{tcolorbox}[colback=verde!8,colframe=verde,title={\faCheck~Lo que ya funciona},fonttitle=\bfseries]
\begin{itemize}[leftmargin=*,itemsep=1pt]
  \item Portada con logo, pregunta al azar y fondo 3D de piezas que cambia en cada visita.
  \item La cartilla, el evento (cuenta regresiva, programa, calendario) y el vistazo con las 40 páginas.
  \item Siete mecanismos en 3D: la balanza interactiva y seis animados; todos en realidad aumentada
        (GLB para Android, USDZ animado para iPhone).
  \item Secciones: GIVIE, Historia, Experiencias, Videos, Galería y Kits.
  \item Galería desde Drive con autor y nombre del modelo; se actualiza sola cada 30 minutos.
  \item Botón «Sube tu diseño» enlazado al formulario.
\end{itemize}
\end{tcolorbox}

\seccion{\faUserShield}{Privacidad y seguridad}
\begin{pendientes}
\casilla & La hoja de respuestas es pública y muestra los \textbf{correos} de quienes suben diseños.
Crear una pestaña solo con fecha, nombre, modelo y archivo; publicarla en la web como CSV; dejar la hoja
completa como \emph{Restringida} y cambiar el enlace \texttt{hoja} en \texttt{datos/config.json}. & GIVIE & \alta\\
\casilla & Moderación: hoy todo lo que se sube al formulario se publica sin revisión. Decidir si se usa
una carpeta aparte solo con los diseños aprobados. & GIVIE & \alta\\
\casilla & Autorización de las familias antes de publicar fotos o videos donde aparezcan niñas y niños
(galería y experiencias). & GIVIE & \alta\\
\casilla & El formulario exige iniciar sesión con Google (por la subida de archivos). Confirmar que no esté
restringido a cuentas de la institución. & GIVIE & \media\\
\end{pendientes}

\seccion{\faCalendarCheck}{Antes del lanzamiento (7 de noviembre)}
\begin{pendientes}
\casilla & \textbf{Revisar toda la información del evento} en la web: fecha, hora, lugar, programa,
responsables y textos, y agregar lo que falte. La \textbf{ubicación probablemente está mal}: confirmar la
dirección exacta del lanzamiento y corregir el botón «Cómo llegar» y el archivo de calendario. & GIVIE & \alta\\
\casilla & Probar la realidad aumentada en un iPhone y en un Android reales, con varios mecanismos. & GIVIE & \alta\\
\casilla & Imprimir el código QR de la web para la proyección y el taller. & GIVIE & \alta\\
\casilla & Reemplazar el video de prueba por los videos reales (idealmente una lista de reproducción de YouTube). & GIVIE / técnico & \media\\
\casilla & Logo de ATA EPE en alta resolución (PNG transparente o SVG). & GIVIE & \media\\
\end{pendientes}

\seccion{\faFile*}{Contenido de las secciones}
\begin{pendientes}
\casilla & \textbf{GIVIE}: miembros, foto y una breve presentación de cada uno. & GIVIE & \media\\
\casilla & \textbf{Experiencias}: escoger fotos y videos de talleres y eventos anteriores (con permisos). & GIVIE & \media\\
\casilla & \textbf{Kits}: qué kits hay, qué traen, precio y cómo conseguirlos. & GIVIE & \media\\
\casilla & \textbf{Historia}: revisar la línea de tiempo; agregar la fundación de la EPE y el nacimiento del kit ATA EPE. & GIVIE & \media\\
\casilla & Revisar las pistas de observación de cada mecanismo con las autoras y autores de la cartilla. & GIVIE & \baja\\
\casilla & Enlazar los códigos QR de la cartilla (publicaciones de \texttt{roboticaescolar.com}). & técnico & \baja\\
\casilla & Después del evento: convertir la sección del lanzamiento en «Así fue», con fotos. & GIVIE / técnico & \baja\\
\end{pendientes}

\seccion{\faCubes}{Modelos 3D}
\begin{pendientes}
\casilla & \textbf{Buscar los manuales de armado} de cada modelo (guías de ensamble como la del cangrejo,
fichas o fotos paso a paso) y pasárselos a Claude, para rehacer los modelos 3D pieza por pieza
y mejorar las aproximaciones. & GIVIE & \alta\\
\casilla & Comparar los seis mecanismos animados con los prototipos reales (fotos de cada uno armado,
de frente y de lado) y corregir el armado. & GIVIE / técnico & \media\\
\casilla & Confirmar los colores reales de cada pieza (en la ilustración, la balanza es verde). & GIVIE & \baja\\
\casilla & Medidas finas: diámetro de los agujeros, tornillos y tuercas reales. & GIVIE & \baja\\
\casilla & Mejorar las figuras (mono, pescador, bailarinas, bote). & técnico & \baja\\
\end{pendientes}

\seccion{\faCogs}{Técnico}
\begin{pendientes}
\casilla & Anexar la web al dominio definitivo cuando haya acceso (es estática: basta copiar la carpeta). & técnico & \media\\
\casilla & Decidir si se dejan públicas las páginas de revisión \texttt{piezas.html} y \texttt{mecanismos.html}. & GIVIE & \baja\\
\casilla & Revisar el rendimiento en celulares de gama baja (fondo 3D y modelos). & técnico & \baja\\
\end{pendientes}

\begin{tcolorbox}[colback=amarillo!15,colframe=amarillo!80!black,title={\faInfoCircle~Cómo se actualiza la web},fonttitle=\bfseries]
\begin{itemize}[leftmargin=*,itemsep=1pt]
  \item \textbf{Videos y galería}: todo se configura en \texttt{datos/config.json}.
  \item \textbf{Galería ya}: en GitHub, \emph{Actions → Actualizar galería → Run workflow}.
  \item \textbf{Quitar un diseño}: sacarlo de la carpeta de Drive; desaparece en la siguiente revisión.
  \item \textbf{Modelos 3D}: \texttt{herramientas/generar\_piezas.py} y \texttt{generar\_mecanismos.py} (ver el README).
\end{itemize}
\end{tcolorbox}

\end{document}
